% AUTHOR: Amlal El Mahrouss
% PURPOSE: DN001.11: Notes for Several Formulas.

\documentclass[11pt, a4paper]{article}
\usepackage{graphicx}
\usepackage{listings}
\usepackage{amsmath,amssymb,amsthm}
\usepackage{xcolor}
\usepackage{hyperref}
\usepackage{mathtools}
\usepackage[margin=0.5in,top=1in,bottom=1in]{geometry}

\title{Fourth Notes.}
\author{Amlal El Mahrouss\\\texttt{amlal@nekernel.org}}
\date{February 23, 2026}

\begin{document}

\maketitle

\begin{center}
	\rule[0.01cm]{17cm}{0.01cm}
\end{center}

\section{DEFINITIONS}

Consider the following formulas: 

\section{DEFINITION OF THE E-VOLATILITY $\operatorname{EVol}_{1}$ OPERATOR:}

Let a definition of $\operatorname{EVol}_{1}$:

\begin{equation}
    \operatorname{EVol}_{1}(n, p) = \int_{n}^{p} \frac{\operatorname{Dec}_{n}(t)}{\operatorname{Diff}_{n}(t)}dt
\end{equation}
Where $\operatorname{Dec}_{n}(x) \in \mathbb{R}$ is the Decremental function of a variable $x$. We assume $\operatorname{Diff}_{n}(x) \in \mathbb{R}$ as the Differential equation of a variable $x$.

\section{DEFINITION OF THE DECREMENTAL FUNCTION $\operatorname{Dec}_{n}(x)$}

Let a function be defined in such manner that its algorithm returns a result in $\mathbb{R}$. Such function is called the $\operatorname{Dec}_{n}(x)$. We define it such that: $ \operatorname{\operatorname{Dec}_{n}(x)} \in \mathbb{R}$

\section{INTRODUCTION}

\end{document}
